\documentclass[11pt]{article}
% Royden report conventions are embedded so this source is standalone.
\usepackage[margin=1in]{geometry}
\usepackage{amsmath,amssymb,amsthm,mathtools}
\usepackage{graphicx,booktabs}
\usepackage[hidelinks]{hyperref}
\setlength{\parindent}{1.25em}
\setlength{\parskip}{0pt}
\pagestyle{plain}
\numberwithin{equation}{section}
\title{Avoiding Cavitation: Problem 5--63}
\author{}
\date{}

\begin{document}
\begin{center}
  {\Large Avoiding Cavitation: Problem 5--63}\par
  \smallskip
  Worked solution
\end{center}

\noindent\textbf{Problem.} Water flows from an open reservoir through a
horizontal pipe. The free surface (3) is $5\,\mathrm m$ above sections (1)
and (2), and the outlet diameter is $D_2=15\,\mathrm{cm}$. Find the minimum
diameter at section (1) that avoids cavitation there.

\medskip
\noindent\textbf{Assumptions.} The flow is steady and incompressible,
losses are negligible, and the reservoir is large enough that $V_3\approx0$.
Both the free surface and outlet are at atmospheric pressure. Use
$\rho=1000\,\mathrm{kg/m^3}$, $g=9.81\,\mathrm{m/s^2}$, and
$p_{\mathrm{atm}}=101.3\,\mathrm{kPa}$. All pressures are absolute.
The temperature is not given; assume $20^\circ\mathrm C$, for which
$p_v\approx2.34\,\mathrm{kPa}$.\footnote{NIST,
\href{https://nvlpubs.nist.gov/nistpubs/jres/80A/jresv80An5-6p775_A1b.pdf}
{\emph{Vapor pressure formulation for water in range 0 to
100 degrees C: A revision}}, Table 3.}

\section{Outlet velocity}
Take $z_1=z_2=0$ and $z_3=5\,\mathrm m$. Bernoulli's equation between
the free surface (3) and outlet (2) gives
\[
 \frac{p_3}{\rho g}+\frac{V_3^2}{2g}+z_3
 =\frac{p_2}{\rho g}+\frac{V_2^2}{2g}+z_2.
\]
Since $p_3=p_2=p_{\mathrm{atm}}$ and $V_3\approx0$, this reduces to
\[
 5=\frac{V_2^2}{2g},
 \qquad
 V_2=\sqrt{2g(5\,\mathrm m)}\approx9.90\,\mathrm{m/s}.
\]

\section{Cavitation limit at section (1)}
Cavitation begins when the local absolute pressure reaches the vapor
pressure. Set $p_1=p_v$ to find the limiting velocity. Bernoulli between
(1) and (2), which are at the same elevation, yields
\[
 \frac{p_v}{\rho g}+\frac{V_{1,\max}^2}{2g}
 =\frac{p_{\mathrm{atm}}}{\rho g}+\frac{V_2^2}{2g}.
\]
Hence
\[
 V_{1,\max}
 =\sqrt{V_2^2+\frac{2(p_{\mathrm{atm}}-p_v)}{\rho}}
 =\sqrt{98.1+\frac{2(101300-2340)}{1000}}\,\mathrm{m/s}
 \approx17.2\,\mathrm{m/s}.
\]

\section{Minimum diameter}
For circular sections, continuity gives
\[
 A_1V_1=A_2V_2,
 \qquad
 \frac{\pi D_1^2}{4}V_1=\frac{\pi D_2^2}{4}V_2.
\]
The smallest allowable diameter corresponds to $V_1=V_{1,\max}$. Thus
\[
 D_{1,\min}
 =D_2\sqrt{\frac{V_2}{V_{1,\max}}}
 \approx0.15\sqrt{\frac{9.90}{17.2}}\,\mathrm m
 \approx\boxed{11.4\,\mathrm{cm}}.
\]
If vapor pressure is neglected ($p_v\approx0$), the same calculation gives
$D_{1,\min}\approx11.3\,\mathrm{cm}$.
A smaller diameter raises the velocity and lowers the pressure at (1).
Equality marks the onset of cavitation; a diameter above the limiting value
keeps $p_1>p_v$ under the stated assumptions.

\end{document}
